\documentclass[10pt,a4paper]{amsart}
\usepackage[margin=19mm]{geometry}
\usepackage{amsmath,amssymb,mathtools}
\usepackage[scr=rsfs,cal=euler]{mathalfa}
\usepackage{xfrac}
\usepackage[unicode,hidelinks,bookmarks=false]{hyperref}
\newcommand{\ie}{i.e.\ }
\newcommand{\cf}{cf.\ }
\newcommand{\resp}{respectively\ }
\newcommand{\Subsection}{\subsection}
\providecommand{\nobreakdash}{\nobreak\hbox{-}}
\setlength{\emergencystretch}{2em}
\multlinegap0pt
\usepackage{amssymb}
\usepackage{mathtools}
\usepackage{bm}
\usepackage[shortlabels]{enumitem}
\newtheorem{theorem}{Theorem}
\newcommand{\norm}[1]{\left\lVert#1\right\rVert}
\title{Gaussian Operator Algebras:\\ Observability, Reconstruction, and Microscopic Limits}
\author{Guangqian Zhao}
\date{Source: arXiv:2606.29292v6 (CC BY 4.0)}
\begin{document}
\maketitle

\noindent\emph{Source and licence.} Gaussian Operator Algebras:\\ Observability, Reconstruction, and Microscopic Limits. Version arXiv:2606.29292v6 (CC BY 4.0), distributed by arXiv under the Creative Commons Attribution 4.0 International licence, http://creativecommons.org/licenses/by/4.0/, as recorded in the arXiv licence metadata for this exact version and shown on the abstract page https://arxiv.org/abs/2606.29292v6. Full licence notice: \texttt{licenses/arxiv-2606.29292-CC-BY-4.0.txt}.\par\smallskip

\noindent\textit{Editorial scope.} Counted unit: the article's theorem thm:arity (source0.tex lines 1073--1084). The article's complete original proof of it is delivered with it (source0.tex lines 1085--1141, one \texttt{proof} environment of 57 lines). No mathematics is written, repaired or re-derived: the statement and the proof are the article's own lines, in the article's own notation, and no retained line is substituted or re-flowed.  Retained uncounted as the source-local context the proof needs: lines 1044--1048.  Nothing was omitted from any retained span, so the package carries no omission record.  No retained line contains a citation, so the wrapper carries no bibliography and no bibliography entry.  No retained line contains a figure environment, an image-inclusion command or drawing code, so no image is delivered and the package declares no asset dependency.  Editorial formatting only, in addition: an AMS \texttt{amsart} preamble that restates the article's own declarations and macros used by the retained lines, namely the environments \texttt{theorem} on separate counters, together with the standard packages those lines need.  The retained environments print in this document as Theorem 1 in order of appearance, whereas the article prints the same items with the numbering of its own full document.  The complete native source of the article as distributed in the arXiv e-print for this exact version is bundled unchanged as \texttt{sources/203923/source0.tex}.
\begin{equation}\label{eq:positivekernel}
 \sum_{r=0}^{m\wedge n}c_{m,n,r}\sigma^{m+n-2r}
 \le\left(\sigma^2+\frac{1+\sigma^2}{t}\right)^{m/2}
     \left(\sigma^2+(1+\sigma^2)t\right)^{n/2}.
\end{equation}

\begin{theorem}[The sharp multi-input radius region]\label{thm:arity}
Assume $H\ne0$, $E\ne0$, $k\ge2$, $\sigma\ge1$, and $\rho_i>0$. Let $\mathcal C_{\rho,2}$ be the kernel space with finite $\mathcal N_{\rho,2}$ norm. Ordered multiplication on the finite core is bounded from $\prod_{i=1}^k\mathcal C_{\rho_i,2}$ into $\mathcal C_{\sigma,2}$ if and only if
\begin{equation}\label{eq:arityregion}
 \sum_{i=1}^k\frac1{1+\rho_i^2}\le\frac1{1+\sigma^2}.
\end{equation}
In that region, its best multilinear constant is one:
\begin{equation}\label{eq:aritybound}
 \mathcal N_{\sigma,2}(K_1\star\cdots\star K_k)
 \le\prod_{i=1}^k\mathcal N_{\rho_i,2}(K_i).
\end{equation}
For equal input radii, the exact condition is $\rho^2\ge k(1+\sigma^2)-1$.
\end{theorem}

\begin{proof}
\emph{Binary sufficiency and absolute convergence.}
For two inputs,~\eqref{eq:arityregion} is equivalent to
\begin{equation}\label{eq:binaryregion}
 \rho^2>\sigma^2,\quad\eta^2>\sigma^2,\quad
 (\rho^2-\sigma^2)(\eta^2-\sigma^2)\ge(1+\sigma^2)^2.
\end{equation}
Indeed, clearing the positive denominators gives the last inequality after adding $\sigma^4$; either of the first two failures makes the reciprocal sum too large. Choose
\[
 \frac{1+\sigma^2}{\rho^2-\sigma^2}\le t
 \le\frac{\eta^2-\sigma^2}{1+\sigma^2}.
\]
Both factors on the right of~\eqref{eq:positivekernel} are then bounded by $\rho^m$ and $\eta^n$. The contraction estimate yields the stronger absolute bound
\begin{equation}\label{eq:absolutecontraction}
 \begin{split}
 &\sum_{m,n\ge0}\sum_{r\le m\wedge n}
 \sigma^{m+n-2r}\sqrt{(m+n-2r)!}\,b_{m,n,r}
       \norm{\Gamma_r(K_m,L_n)}\\
 &\hspace{25mm}\le\mathcal N_{\rho,2}(K)\mathcal N_{\eta,2}(L).
 \end{split}
\end{equation}
All contractions can therefore be grouped by output degree. Finite kernels are dense in the input spaces, so the bound supplies the unique continuous extension.

\emph{Ordered multi-input multiplication.}
Set $c_i=(1+\rho_i^2)^{-1}$. Choose an arbitrary binary bracketing preserving input order. For a node containing the consecutive input set $A$, assign the radius
\[
 r_A=\left(\left(\sum_{i\in A}c_i\right)^{-1}-1\right)^{1/2}.
\]
The root radius is at least $\sigma$, and every intermediate radius is at least the root radius. For each binary node, the reciprocal weights of its two children add to that of their parent, so the binary estimate holds at equality in its radius condition. Iterating with constant one proves~\eqref{eq:aritybound}; monotonicity in the output radius gives the prescribed $\sigma$. Associativity on the finite core and continuity make the result independent of bracketing. Taking every input to be the same constant rank-one projection of Hilbert--Schmidt norm one proves that the constant cannot be smaller than one.

\emph{Necessity.}
Fix a unit source vector $h$ and a rank-one projection $P$. For $z\ge0$, consider the coherent kernel
\[
 \kappa(z)_m=\frac{z^m}{m!}h^{\otimes m}\otimes P,
 \qquad\mathcal A(u)=\sum_{m\ge0}\frac{u^m}{\sqrt{m!}}.
\]
Direct contraction gives
\[
 \kappa(z_1)\star\cdots\star\kappa(z_k)
 =\exp\left(\sum_{i<j}z_iz_j\right)\kappa\left(\sum_i z_i\right),
 \qquad\mathcal N_{\rho,2}(\kappa(z))=\mathcal A(\rho z).
\]
These tests are justified even when boundedness is initially assumed only on the finite core: truncate each input, and pass to the limit in the nonnegative coefficient sums.

We have $\mathcal A(u)^2\ge e^{u^2}$. For $0<q<1$, Cauchy--Schwarz gives
\[
 \mathcal A(u)\le(1-q^2)^{-1/2}\exp\big(u^2/(2q^2)\big).
\]
Letting $u\to\infty$ and then $q\uparrow1$ shows
$\log\mathcal A(u)/u^2\to1/2$.
If multiplication has any finite bound, substitute $z_i=ta_i$, take logarithms, divide by $t^2$, and let $t\to\infty$. One obtains
\[
 \sum_i(1+\rho_i^2)a_i^2-(1+\sigma^2)\left(\sum_i a_i\right)^2\ge0
 \qquad(a_i\ge0).
\]
Taking $a_i=(1+\rho_i^2)^{-1}$ gives~\eqref{eq:arityregion}.
\end{proof}

\end{document}
